\textbf{From symbolic planning to LLM/VLM planners.} Robotic planning
has long been structured around a Sense-Plan-Act
pipeline~\cite{garrett2021integrated}, first realized through
hierarchical and symbolic planners that search a discrete action space
for a provably goal-reaching sequence~\cite{kaelbling2011hierarchical}.
These planners give strong guarantees but require a hand-specified
world model. Foundation models relaxed that requirement. They ground
natural-language instructions in robotic
affordances~\cite{ahn2022saycan}, act as zero-shot
planners~\cite{huang2022zeroshot}, and generate motion plans directly
from instructions~\cite{lin2023text2motion}. On top of these
capabilities, a range of reasoning strategies has been proposed:
reactive single-step prompting, interleaved reasoning and acting
(ReAct)~\cite{yao2022react}, iterative
self-critique~\cite{madaan2023selfrefine, zhou2024isrllm}, tree
search~\cite{yao2023tot}, and chain-of-thought with self-consistency
voting~\cite{wei2022cot, wang2022selfconsistency}. Recent surveys
catalogue this space~\cite{kim2024survey, wang2025robotics,
huang2024planningsurvey, berti2025emergent, billard2025roadmap,
favaliRAS}. Our contribution is orthogonal: we propose no new
reasoning strategy, and our evaluation holds the strategy fixed. We
address a failure axis that sits underneath all of them: how a chosen
action's numeric parameters are computed. That axis persists whatever
strategy produced the action.

\textbf{Why Sense-Plan-Act with LLM/VLM planners over a VLA.} An
alternative architecture collapses perception, reasoning, and control
into one vision-language-action policy~\cite{brohan2023rt2,
driess2023palme, geminiteam2025robotics15, kim2024openvla}. We argue
for the decoupled architecture on two grounds. The first is data. A VLA
must be trained or fine-tuned on paired vision-action data, which is
scarce, embodiment-specific, and costly to collect. A decoupled
pipeline instead plans in text, or text conditioned on an image, and
emits a sequence of parameterized skills. The model's pretrained
knowledge, spatial reasoning, arithmetic, common-sense physics, and
ordering logic, becomes available zero-shot, through code, without a
single task-specific demonstration~\cite{liang2023codeaspolicies,
huang2023voxposer, huang2022innermonologue, singh2023progprompt}. The
second ground is diagnosability. When an end-to-end policy fails, it is
hard to tell a misperception from a reasoning error or a motor-control
error~\cite{favaliRAS}. A decoupled pipeline exposes its intermediate
plan, so every failure can be traced to a specific stage and a specific
model output. The failure analysis of this paper
(Section~\ref{sec:failures}) would not have been possible in an
end-to-end system; only its symptoms would have been visible. The cost
of decoupling is that the model must produce numeric parameters it was
never trained to get right. \rev{The harness proposed in this paper
addresses precisely this weakness.}

\textbf{Simulation, sim-to-real, and the closest prior work.} The
sim-to-real gap is a long-studied concern in robot learning, primarily
for learned low-level control~\cite{zhao2020simtoreal}. Physics
simulators such as Gazebo~\cite{koenig2004gazebo},
MuJoCo~\cite{todorov2012mujoco}, and
CoppeliaSim~\cite{rohmer2013coppeliasim} narrow that gap by modeling
contact dynamics and actuation limits~\cite{collins2021physics}. The
gap we document is different in kind. It concerns whether a
\emph{model-computed} geometric quantity is trustworthy at all, a
question a symbolic benchmark never poses, because it never confronts
the model's number with a contradicting sensor. The closest prior work
is Favali et al.~\cite{favaliRAS}, which introduces a task-complexity
taxonomy and the TS/TSR/AETS metrics for benchmarking LLM-based robotic
planning across six reasoning strategies, evaluated in a symbolic
simulator with exact ground truth. We adopt the same evaluation
philosophy and the same metric family, but on real hardware, where
ground truth for safety and success is not automatically available.
